\documentclass[11pt]{article}
\usepackage{fontspec}
\usepackage{amsmath,amssymb,amsthm,mathtools}
\usepackage[margin=1.1in]{geometry}
\usepackage{microtype}
\usepackage{parskip}
\usepackage{lmodern}
\usepackage{xcolor}
\usepackage{adjustbox}
\usepackage[colorlinks=true,linkcolor=blue!50!black,urlcolor=blue!50!black]{hyperref}
\providecommand{\tightlist}{}
\setlength{\emergencystretch}{3.5em}
\title{The delta orientation is the unique entropy minimizer for self-calibrated ±1 Walsh tilts on the Boolean cube}
\author{Felix Robles Elvira (ORCID: 0009-0009-2017-4394; independent researcher)}
\date{\small preprint, not peer reviewed, version 2026-07-02.  This paper is mathematically self-contained: no result below depends on any other paper in this repository or elsewhere.  The range $n \le 5$ of the main theorem is computer-assisted; Section 8 states the exact computational claims, their certificates, and their trust base, and the repository ships the complete certification code and canonical outputs (receipts \texttt{w1}–\texttt{w5}).}
\begin{document}
\maketitle

\section*{Abstract}
For each choice of signs $\epsilon_a \in \{\pm 1\}$, one per nonzero Walsh character $\chi_a$ on the Boolean cube $G = \{\pm1\}^n$ ($N = 2^n$ points), consider the exponential-family law $P_\epsilon(s) \propto \exp\big(\sum_{a \ne 0} h_a \epsilon_a \chi_a(s)\big)$ whose parameters are \emph{self-calibrated} by the fixed-point conditions $\mathbb{E}_{P_\epsilon}[\epsilon_a \chi_a] = e^{-h_a}$ for every $a \ne 0$: each tilted correlation must equal the exponential of minus its own tilt parameter.  We prove that this system has exactly one solution for every sign choice, and that the relative entropy $D(P_\epsilon \Vert U)$ from the uniform law is minimized, strictly and uniquely up to the natural $N$-element symmetry orbit, by the \textbf{delta orientations} $\epsilon_a = -\chi_a(s_\star)$ — the sign patterns whose calibrated law is uniform on $N-1$ points and nearly extinguishes the single point $s_\star$.  The proof is elementary and quantitative.  Its core is a \emph{deep-dip trichotomy}: any non-delta calibrated law with $D \le 1/60$ must extinguish at least three points of the cube to depth $e^{-5}$, forcing $N \cdot D \ge 3\,\psi(e^{-5}) > 2.878716$ where $\psi(x) = x\log x - x + 1$, while the delta value satisfies $N \cdot D_\delta < N/(N-1)$.  For $n \ge 6$ this closes the theorem analytically; for $n \le 5$ we verify it by certified computation (exhaustive over all $2^{N-1}$ orientations for $n \le 4$; via a provably complete $176$-orbit symmetry transversal of the $2^{31}$ orientations at $n = 5$), with a posteriori Newton–Kantorovich error certificates throughout.  The minimizing value is not a photo-finish: the best non-delta family we exhibit satisfies $N \cdot D \to 4\log 4 = 5.545\ldots$, and we conjecture this is the sharp asymptotic runner-up constant.  The self-calibration fixed point appears to be new; its nearest relatives are the extremal theory of Littlewood ($\pm1$-coefficient) polynomials and the self-consistent field equations of statistical mechanics, and we state our literature position explicitly in Section 1.4.

\textbf{MSC 2020:} 94A17 (primary); 60E15, 42C10, 05D40, 65G20 (secondary).

\textbf{Keywords:} Boolean cube; Walsh--Fourier analysis; relative entropy; exponential family; self-calibration; computer-assisted proof; Newton--Kantorovich certificate.

\section*{1. Introduction}
\subsection*{1.1 The objects}
Let $n \ge 2$, let $G = \{-1,+1\}^n$ be the Boolean cube with $N = 2^n$ points, and let $U$ denote the uniform probability law on $G$.  We identify $G$ with $\mathbb{F}_2^n$ (coordinatewise, $+1 \leftrightarrow 0$ and $-1 \leftrightarrow 1$), write $s + t$ for the group operation (coordinatewise multiplication in the $\pm1$ picture), and index the Walsh characters by masks $a \in \mathbb{F}_2^n$:

\[
\chi_a(s) = \prod_{i \,:\, a_i = 1} s_i, \qquad
\chi_0 \equiv 1, \qquad
\chi_a \chi_b = \chi_{a+b}.
\]
The $N$ characters form an orthonormal basis of $L^2(G, U)$; for $f : G \to \mathbb{R}$ we write $\widehat f(a) = \mathbb{E}_U[f \chi_a]$, so that $f = \sum_a \widehat f(a) \chi_a$ and Parseval's identity $\mathbb{E}_U[f g] = \sum_a \widehat f(a) \widehat g(a)$ holds.

An \textbf{orientation} is a sign vector $\epsilon = (\epsilon_a)_{a \ne 0} \in \{\pm1\}^{N-1}$, one sign per nonzero mask.  There are $2^{N-1}$ orientations.

\begin{quote}
\textbf{Definition 1.1 (calibrated law).}  A probability law $P$ on $G$ with full support is \emph{calibrated for the orientation $\epsilon$} if there exist reals $(h_a)_{a \ne 0}$ such that
\[
P(s) = \frac{\exp\big(\sum_{a\ne0} h_a \epsilon_a \chi_a(s)\big)}
{\sum_{r \in G} \exp\big(\sum_{a\ne0} h_a \epsilon_a \chi_a(r)\big)}
\qquad \text{and} \qquad
\mathbb{E}_P[\epsilon_a \chi_a] = e^{-h_a} \quad (a \ne 0).
\]
\end{quote}
The defining conditions are a fixed point: the parameter $h_a$ shapes $P$, and the correlation that $P$ produces on the signed character $\epsilon_a\chi_a$ is required to equal $e^{-h_a}$, a strictly decreasing function of that same parameter.  Two structural consequences are immediate and are used throughout.  First, since $e^{-h_a} > 0$, a calibrated law has \textbf{nonzero correlation of prescribed sign with every one of the $N-1$ characters}; in particular it is never uniform.  Second, writing $X = N P$ for the density of $P$ with respect to $U$ and $x_a = \widehat X(a) = \mathbb{E}_P[\chi_a]$, the conditions say exactly

\[
x_a = \epsilon_a e^{-h_a}, \qquad
\widehat{\log X}(a) = \epsilon_a h_a \qquad (a \ne 0),
\]
because $\log X = \sum_{a \ne 0} h_a\epsilon_a\chi_a + \text{const}$. Note that $h_a > 0$ automatically: $e^{-h_a} = |x_a| = |\mathbb{E}_P[\chi_a]| < 1$ for any full-support law, since $\chi_a$ is a nonconstant $\pm1$-valued function. Thus the Walsh coefficients of $X$ and of $\log X$ carry the same signs, with magnitudes linked by $|x_a| = e^{-|\widehat{\log X}(a)|}$. The orientation is recoverable from the law: $\epsilon_a = \operatorname{sgn}(x_a)$.

Theorem 3.1 below proves that \textbf{every orientation has exactly one calibrated law}, denoted $P_\epsilon$, as the unique minimizer of an explicit smooth, strictly convex, coercive function on $\mathbb{R}^{N-1}$.  The quantity of interest is the relative entropy (Kullback–Leibler divergence) from uniform,

\[
\widehat m(\epsilon) \;=\; D(P_\epsilon \Vert U)
\;=\; \mathbb{E}_U[X \log X] \;\ge\; 0 ,
\]
which is strictly positive for every orientation (a calibrated law is never uniform).  Which orientation makes it smallest?

\subsection*{1.2 The delta orientations and the main theorem}
Fix $s_\star \in G$ and take $\epsilon_a^\star = -\chi_a(s_\star)$ for every $a \ne 0$: every sign is chosen to \emph{disagree} with the point $s_\star$.  We call this the \textbf{delta orientation at $s_\star$}; there are exactly $N$ of them, one per point, and they form a single orbit under the translation symmetry of the problem (Lemma 3.2).  The name comes from the raw sign field itself: $\sum_{a\ne0} \epsilon^\star_a\chi_a(s) = 1 - N\mathbf{1}_{s = s_\star}$, a negative delta spike.  Proposition 4.1 computes the calibrated law of a delta orientation in closed form: it is the two-level law

\[
X_\delta(s_\star) = 1 - (N-1)u_\star, \qquad
X_\delta(s) = 1 + u_\star \quad (s \ne s_\star),
\]
where $u_\star \in (0, \tfrac1{N-1})$ is the unique positive root of $u^{N+1} + u^N + (N-1)u - 1 = 0$ — an almost-uniform law with one almost-extinguished point, whose gap $D_\delta := \widehat m(\epsilon^\star)$ satisfies $D_\delta < \frac{1}{N-1}$ (Lemma 4.2) and in fact $N D_\delta \to 1$.

\begin{quote}
\textbf{Theorem 1.2 (main theorem).}  For every $n \ge 2$ and every orientation $\epsilon$,
\[
\widehat m(\epsilon) \;\ge\; \widehat m(\epsilon^\star),
\]
with equality if and only if $\epsilon$ is one of the $N$ delta orientations.
\end{quote}
The theorem is proved analytically for all $n \ge 6$ (Sections 5–7) and by certified computation for $2 \le n \le 5$ (Section 8).  The analytic part yields the quantitative floor

\begin{quote}
\textbf{Corollary 1.3.}  For $n \ge 6$ and every non-delta orientation $\epsilon$: $\; N \cdot \widehat m(\epsilon) \ge \min\!\big(N/60,\; 2.878716\big)$, while $N \cdot \widehat m(\epsilon^\star) < N/(N-1) \le 64/63$.
\end{quote}
\subsection*{1.3 The proof in one paragraph}
Suppose $\epsilon$ is not a delta orientation and $D = \widehat m(\epsilon) \le 1/60$; we must show this is very expensive anyway. Pinsker's inequality forces every correlation $|x_a| \le \sqrt{2D}$, hence every parameter $h_a \ge \frac12\log\frac{1}{2D} \ge \frac12\log 30$: \emph{small entropy gap forces uniformly large parameters.}  Since $\widehat{\log X}(a) = \epsilon_a h_a$, the function $\log X$ has huge Fourier coefficients in every direction, and an $\ell^1$ bookkeeping of $\log X$ over the cube (Section 5) shows these coefficients are reproduced, up to a per-mask error $\varepsilon_G \le 2.5\sqrt{2D}$, by the \emph{dip transform} $\Phi(a) = \sum_s \beta_s \chi_a(s)$ of the deeply suppressed points ($X(s) \le \frac12$, depth $\beta_s = -\log X(s)$).  Consequently $|\Phi(a)| \ge N(h_{\min} - \varepsilon_G) > 1.24416\,N$ at every mask, and the signs of $\Phi$ determine the orientation: $\epsilon_a = -\operatorname{sgn} \Phi(a)$.  Now split the dips at depth $5$.  The total depth of shallow dips is less than $0.543149\,N$ (their number is controlled by the entropy budget), so: with \emph{zero} deep dips, $\Phi$ cannot clear its floor; with \emph{one}, that dip dominates the configuration and the sign readout says $\epsilon$ is a delta orientation; with \emph{two}, the floor on the anti-symmetric mask class forces a depth separation $>0.701014\,N$, which again makes the deeper dip dominant — a contradiction in all three cases.  Hence at least \textbf{three} points are suppressed below $e^{-5}$, and since $D = \mathbb{E}_U[\psi(X)]$ with $\psi \ge 0$, these alone cost $N D \ge 3\psi(e^{-5}) > 2.878716$. The delta costs less than $N/(N-1) \le 64/63$; and if instead $D > 1/60$ then $D > 1/60 > 1/63 \ge 1/(N-1) > D_\delta$ directly. Either way the delta wins strictly, for every $n \ge 6$.

We emphasize the shape of the argument: the theorem is \emph{not} won by locating the best non-delta competitor.  The cheapest non-delta family we know (Section 9) costs $N D \to 4\log4 = 5.545\ldots$, and proving that sharp constant remains open; but the theorem only requires beating $N D_\delta < 64/63$, and "three extinguished points versus one" clears that bar with margin $2.8\times$.

\subsection*{1.4 Relation to the literature}
We have not found the self-calibration fixed point of Definition 1.1 in the literature, and we state this as a position, not a claim of certainty.  Its ingredients are all classical.  (i) The sign fields $T_\epsilon = \sum_{a\ne0}\epsilon_a\chi_a$ are the Boolean-cube analogue of Littlewood polynomials (all coefficients $\pm1$), whose extremal theory — flatness, merit factors — is a long-standing subject [3]; our problem differs in that the field is passed through a nonlinear, self-referential Gibbs calibration before being measured. (ii) Exponential families on the cube with Walsh sufficient statistics are the standard log-linear/graphical models [4]; there the parameters are free or fit to external data, whereas here they are tied to the model's own correlations.  (iii) Self-consistent equations linking a field to its own expectation are the daily bread of mean-field statistical mechanics [6], but the standard link is $m = \tanh(h)$; the link $m = e^{-h}$, imposed \emph{coefficientwise across an entire Walsh spectrum}, changes the character of the problem entirely (for instance, it forbids any vanishing correlation).  (iv) The analysis lives in the standard framework of Boolean Fourier analysis [2] and uses only Pinsker's inequality [1] and elementary convexity.

\subsection*{1.5 Organization}
Section 2 fixes notation and records the two standard inequalities we use.  Section 3 proves existence, uniqueness, and the symmetry covariances of calibrated laws.  Section 4 computes the delta laws in closed form and proves $D_\delta < 1/(N-1)$.  Section 5 establishes the structure theory of low-entropy calibrated laws: the parameter floor, the dip/spike/bulk decomposition with its bookkeeping, the spectral floor and sign readout, and the dominance criterion.  Section 6 proves the deep-dip trichotomy.  Section 7 assembles the main theorem for $n \ge 6$.  Section 8 specifies and certifies the computation for $n \le 5$.  Section 9 presents the extremal family and the sharp-constant conjecture.  Section 10 documents the code.  All decimal constants in Sections 5–7 are rounded in the \emph{safe} direction (lower bounds down, upper bounds up).

\section*{2. Preliminaries}
Throughout, $\log$ is the natural logarithm, $n \ge 2$, $N = 2^n$, and $X = NP$ denotes the density of a probability law $P$ with respect to $U$, so $\mathbb{E}_U[X] = 1$.  The relative entropy is $D(P \Vert U) = \mathbb{E}_U[X\log X]$.  We use two standard facts.

\begin{quote}
\textbf{Lemma 2.1 (Pinsker; see [1], Lemma 11.6.1).}  For any probability law $P$ on $G$,
\[
\mathbb{E}_U\,|X - 1| \;=\; \sum_{s}\,\big|P(s) - \tfrac1N\big|
\;\le\; \sqrt{2\,D(P\Vert U)} .
\]
\end{quote}
(The middle quantity is twice the total variation distance; Pinsker's inequality in the form $\mathrm{TV} \le \sqrt{D/2}$ gives the stated bound.)  Two consequences used constantly: every Walsh coefficient obeys $|x_a| = |\mathbb{E}_U[(X-1)\chi_a]| \le \mathbb{E}_U|X-1| \le \sqrt{2D}$; and, since $\mathbb{E}_U[X-1] = 0$, the positive and negative parts match, $\mathbb{E}_U(X-1)_+ = \mathbb{E}_U(1-X)_+ = \tfrac12\mathbb{E}_U|X-1| \le \tfrac12\sqrt{2D}$.

\begin{quote}
\textbf{Lemma 2.2 (entropy integrand).}  Let $\psi(x) = x\log x - x + 1$ for $x > 0$, $\psi(0) = 1$.  Then $\psi \ge 0$, $\psi$ is strictly decreasing on $[0,1]$ and strictly increasing on $[1,\infty)$, and for any law $P$,
\[
D(P \Vert U) \;=\; \mathbb{E}_U[\psi(X)]
\;=\; \frac1N \sum_{s} \psi(X(s)).
\]
\end{quote}
\emph{Proof.}  $\psi'(x) = \log x$, which is negative on $(0,1)$ and positive on $(1,\infty)$; $\psi(1) = 0$, so $\psi \ge 0$ and the monotonicity is as stated; $\psi(x) \to 1$ as $x \downarrow 0$.  The identity follows from $\mathbb{E}_U[X] = 1$: $\mathbb{E}_U[\psi(X)] = \mathbb{E}_U[X\log X] - 1 + 1 = D$.  $\square$

We will also need two elementary calculus bounds, proved where used (Lemma 5.2): $|\log x| \le 2|x-1|$ on $[\tfrac12, 2]$, and specific values of $\psi$.

\section*{3. Existence, uniqueness, and symmetry of calibrated laws}
For $\ell \in \mathbb{R}^{N-1}$ (indexed by nonzero masks) let

\[
F(\ell) \;=\; \log \mathbb{E}_U \exp\Big(\sum_{a\ne0} \ell_a
\chi_a\Big)
\]
be the log-partition function and $P_\ell \propto \exp(\sum_a \ell_a\chi_a)$ the associated Gibbs law, so that $\partial F/\partial \ell_a = \mathbb{E}_{P_\ell}[\chi_a] =: x_a(\ell)$ and $\nabla^2 F(\ell) = \operatorname{Cov}_{P_\ell}(\chi)$, the covariance matrix of the character vector.

\begin{quote}
\textbf{Theorem 3.1 (existence and uniqueness).}  For every orientation $\epsilon$ the function
\[
G_\epsilon(\ell) \;=\; F(\ell) \;+\; \sum_{a \ne 0}
e^{-\epsilon_a \ell_a}
\]
is smooth, strictly convex, and coercive on $\mathbb{R}^{N-1}$; its unique global minimizer $\ell^\ast$ satisfies the calibration conditions of Definition 1.1 with $h_a = \epsilon_a\ell^\ast_a > 0$, and conversely every calibrated law for $\epsilon$ arises this way. In particular each orientation has exactly one calibrated law $P_\epsilon$.
\end{quote}
\emph{Proof.}  \emph{Smoothness and strict convexity.}  $F$ is smooth with $\nabla^2 F = \operatorname{Cov}_{P_\ell}(\chi) \succeq 0$; each barrier term $e^{-\epsilon_a\ell_a}$ is smooth and convex with second derivative $e^{-\epsilon_a\ell_a} > 0$ in the coordinate $\ell_a$, so $\nabla^2 G_\epsilon \succeq \operatorname{diag}(e^{-\epsilon_a \ell_a}) \succ 0$: $G_\epsilon$ is strictly convex.

\emph{Coercivity.}  We first show $G_\epsilon(t v) \to \infty$ as $t \to \infty$ for every unit vector $v \ne 0$.  The function $s \mapsto \langle v, \chi(s)\rangle = \sum_a v_a\chi_a(s)$ has zero mean over $G$ and is not identically zero (the characters are linearly independent), so $\mu(v) := \max_s \langle v, \chi(s)\rangle > 0$. Bounding the expectation defining $F$ from below by the single term at an attaining $s$,

\[
F(tv) \;\ge\; t\,\mu(v) - \log N \;\longrightarrow\; \infty ,
\]
and the barrier terms are nonnegative, so $G_\epsilon(tv) \to \infty$.  Now let $\ell_k$ be any sequence with $\Vert\ell_k\Vert \to \infty$ and suppose $G_\epsilon(\ell_k) \le C$.  Passing to a subsequence, $v_k := \ell_k/\Vert\ell_k\Vert \to v$, a unit vector. Fix $t > 0$.  For all large $k$, $t \le \Vert\ell_k\Vert$, so convexity along the segment from $0$ to $\ell_k$ gives $G_\epsilon(t v_k) \le \max\{G_\epsilon(0), G_\epsilon(\ell_k)\} \le \max\{G_\epsilon(0), C\}$.  A finite convex function on $\mathbb{R}^{N-1}$ is continuous, so letting $k \to \infty$: $G_\epsilon(tv) \le \max\{G_\epsilon(0), C\}$ for every $t$ — contradicting $G_\epsilon(tv) \to \infty$.  Hence $G_\epsilon$ is coercive and attains a unique global minimum $\ell^\ast$.

\emph{The critical equations are the calibration.}  At any critical point,

\[
0 = \frac{\partial G_\epsilon}{\partial \ell_a}
= x_a(\ell) - \epsilon_a e^{-\epsilon_a\ell_a}
\quad\Longleftrightarrow\quad
\mathbb{E}_{P_\ell}[\epsilon_a\chi_a] = e^{-\epsilon_a \ell_a} ,
\]
which is Definition 1.1 with $h_a = \epsilon_a\ell_a$; and $h_a > 0$ follows because $e^{-h_a} = |\mathbb{E}_{P}[\chi_a]| < 1$ for any law with full support (the Gibbs law has full support, and $|\chi_a| = 1$ with $\chi_a$ nonconstant).  Conversely a calibrated law is a Gibbs law $P_\ell$ with $\ell_a = h_a\epsilon_a$ satisfying the same first-order equations, i.e. a critical point of the strictly convex $G_\epsilon$, hence \emph{the} minimizer.  $\square$

The problem has an exact symmetry group, which we record next; it is used for the delta computation (Section 4) and for the $n = 5$ orbit reduction (Section 8).

\begin{quote}
\textbf{Lemma 3.2 (translation covariance).}  For $t \in G$ define $(\tau_t\epsilon)_a = \epsilon_a\chi_a(t)$.  Then $P_{\tau_t\epsilon}(s) = P_\epsilon(s + t)$ and $\widehat m(\tau_t\epsilon) = \widehat m(\epsilon)$.  The $N$ delta orientations form a single translation orbit: $\tau_t(\epsilon^\star\text{ at }s_\star) = \epsilon^\star\text{ at } s_\star + t$.
\end{quote}
\emph{Proof.}  Given $\ell$, set $\ell'_a = \ell_a \chi_a(t)$.  Then $\sum_a \ell'_a \chi_a(s) = \sum_a \ell_a \chi_a(s+t)$, so $F(\ell') = F(\ell)$ (substituting $s \mapsto s + t$ in the defining expectation, a bijection preserving $U$) and $P_{\ell'}(s) = P_\ell(s+t)$.  Also $(\tau_t\epsilon)_a \ell'_a = \epsilon_a\ell_a$, so the barrier sums agree and $G_{\tau_t\epsilon}(\ell') = G_\epsilon(\ell)$.  The map $\ell \mapsto \ell'$ is a linear bijection, so it carries the unique minimizer to the unique minimizer; the laws are translates of one another, and relative entropy to the (translation-invariant) uniform law is unchanged.  For the delta orbit: $(\tau_t \epsilon^\star)_a = -\chi_a(s_\star)\chi_a(t) = -\chi_a(s_\star + t)$.  Since the map $s_\star \mapsto (-\chi_a(s_\star))_{a\ne0}$ is injective (characters separate points), the $N$ delta orientations are distinct and form one orbit.  $\square$

\begin{quote}
\textbf{Lemma 3.3 ($GL(n,2)$ covariance).}  For $M \in GL(n,2)$ define the relabeled orientation $M\cdot\epsilon$ by $(M\cdot\epsilon)_b = \epsilon_{M^{-\mathsf T} b}$, equivalently $(M\cdot\epsilon)_{M^{\mathsf T} a} = \epsilon_a$.  Then $P_{M\cdot\epsilon}$ is the image of $P_\epsilon$ under the $U$-preserving bijection $s \mapsto M^{-1}s$, and $\widehat m(M\cdot\epsilon) = \widehat m(\epsilon)$.  Moreover $M\cdot(\text{delta at } s_\star) = \text{delta at } M^{-1}s_\star$: the family of delta orientations is permuted among itself.
\end{quote}
\emph{Proof.}  The map $s \mapsto Ms$ (matrix action on $\mathbb{F}_2^n$-coordinates) is a bijection of $G$ preserving $U$, and $\chi_a(Ms) = (-1)^{\langle a, Ms\rangle} = (-1)^{\langle M^{\mathsf T} a, s\rangle} = \chi_{M^{\mathsf T}a}(s)$. Given $\ell$, define $\ell'$ by $\ell'_{M^{\mathsf T}a} = \ell_a$. Then $\sum_a \ell_a \chi_a(Ms) = \sum_b \ell'_b \chi_b(s)$, so $F(\ell') = F(\ell)$ (substitute $s \mapsto Ms$ inside the defining expectation) and $P_{\ell'}$ is the image of $P_\ell$ under $s \mapsto M^{-1}s$; the barrier sums agree under the reindexing, since $(M\cdot\epsilon)_{M^{\mathsf T}a}\,\ell'_{M^{\mathsf T}a} = \epsilon_a \ell_a$, so $G_{M\cdot\epsilon}(\ell') = G_\epsilon(\ell)$, and the linear bijection $\ell \mapsto \ell'$ carries the unique minimizer to the unique minimizer.  Relative entropy to $U$ is invariant under a $U$-preserving bijection.  For the delta family, using $\langle M^{-\mathsf T}b, s_\star\rangle = \langle b, M^{-1}s_\star\rangle$:

\[
(M\cdot\epsilon^\star)_b \;=\; \epsilon^\star_{M^{-\mathsf T}b}
\;=\; -\chi_{M^{-\mathsf T}b}(s_\star)
\;=\; -\chi_b(M^{-1}s_\star),
\]
the delta orientation at $M^{-1}s_\star$.  $\square$

Together, Lemmas 3.2 and 3.3 say that $\widehat m$ is constant on the orbits of the group $\Gamma_n = \langle \text{translations}, GL(n,2)\rangle$ acting on orientations, and that the $N$ delta orientations form exactly one $\Gamma_n$-orbit (translations already act transitively on them, and $GL(n,2)$ maps the family to itself).

\section*{4. The delta laws in closed form}
\begin{quote}
\textbf{Proposition 4.1 (delta law).}  Let $\epsilon^\star$ be the delta orientation at $s_\star$.  Its calibrated law is the two-level law
\[
X_\delta(s_\star) = 1 - (N-1)u_\star, \qquad
X_\delta(s) = 1 + u_\star \quad (s \ne s_\star),
\]
where $u_\star$ is the unique positive root of
\[
p(u) \;=\; u^{N+1} + u^N + (N-1)u - 1 \;=\; 0,
\]
and $u_\star \in (0, \tfrac1{N-1})$.  Consequently
\[
D_\delta \;=\; \frac1N\Big[A\log A + (N-1)B\log B\Big],
\qquad A = 1-(N-1)u_\star,\; B = 1+u_\star .
\]
\end{quote}
\emph{Proof.}  By Lemma 3.2 it suffices to treat $s_\star = 0$ (the identity of $G$), where $\chi_a(s_\star) = 1$ and $\epsilon^\star \equiv -1$.  For any $M \in GL(n,2)$, the reindexing of Lemma 3.3 fixes the constant orientation $\epsilon \equiv -1$; hence, by uniqueness of the minimizer, $\ell^\ast_{M^{\mathsf T}a} = \ell^\ast_a$ for all $a \ne 0$ and all $M$.  Since $GL(n,2)$ acts transitively on nonzero masks, all coordinates of $\ell^\ast$ are equal: $\ell^\ast_a = -h$ for some $h > 0$ (Theorem 3.1 gives $h_a = \epsilon_a\ell_a = h > 0$), and $u := e^{-h} \in (0,1)$.

The density is then explicit.  With $x_a = \epsilon_a u = -u$ for all $a \ne 0$ and the Walsh inversion $\sum_{a \ne 0}\chi_a(s) = N\mathbf{1}_{s=0} - 1$:

\[
X(s) = 1 + \sum_{a\ne0} x_a \chi_a(s)
= 1 - u\big(N\mathbf{1}_{s=0} - 1\big)
= \begin{cases}
1 - (N-1)u, & s = 0,\\
1 + u, & s \ne 0 .
\end{cases}
\]
Positivity of the Gibbs law forces $1 - (N-1)u > 0$, i.e. $u < \frac1{N-1}$.  It remains to impose the calibration on the $\log$-side.  With $A = 1-(N-1)u$, $B = 1+u$,

\[
\widehat{\log X}(a)
= \frac1N\Big[\log A \cdot \chi_a(0) + \log B \sum_{s \ne 0}
\chi_a(s)\Big]
= \frac1N\big(\log A - \log B\big)
\qquad (a \ne 0),
\]
using $\sum_{s\ne0}\chi_a(s) = -1$.  The calibration requires $\widehat{\log X}(a) = \epsilon_a h_a = -h = \log u$, i.e.

\[
\frac{1 - (N-1)u}{1+u} \;=\; u^N ,
\]
which after clearing denominators is exactly $p(u) = 0$.  Finally $p$ is strictly increasing on $[0,\infty)$ (its derivative is positive), $p(0) = -1 < 0$, and $p(\tfrac1{N-1}) = (\tfrac1{N-1})^{N}(1 + \tfrac1{N-1}) > 0$, so $p$ has a unique positive root and it lies in $(0, \tfrac1{N-1})$.  Conversely the two-level law with $u = u_\star$ \emph{is} calibrated (all the above read backwards), so by uniqueness (Theorem 3.1) it is the calibrated law of $\epsilon^\star$.  The formula for $D_\delta$ is the definition of $D(P\Vert U)$ evaluated on the two levels.  $\square$

\begin{quote}
\textbf{Lemma 4.2 (elementary delta bound).}  For every $n \ge 2$,
\[
0 \;<\; D_\delta \;<\; \frac{1}{N-1} .
\]
\end{quote}
\emph{Proof.}  Positivity holds because the calibrated law is not uniform (Section 1.1).  For the upper bound: $A \in (0,1)$, so $A\log A \le 0$, and hence

\[
D_\delta \le \frac{N-1}{N}\,(1+u_\star)\log(1+u_\star)
\le \frac{N-1}{N}\,u_\star(1+u_\star),
\]
using $\log(1+x) \le x$.  Now $u_\star < \frac1{N-1}$ gives $1 + u_\star < \frac{N}{N-1}$ and therefore

\[
D_\delta \;<\; \frac{N-1}{N}\cdot\frac{1}{N-1}\cdot\frac{N}{N-1}
\;=\; \frac{1}{N-1}. \qquad \square
\]
Two remarks.  First, the bound is asymptotically tight: from $p(u_\star) = 0$ one gets $u_\star = \frac1{N-1}(1 - O(N^{-N}))$ and $N D_\delta \to 1$; we never need this.  Second, by Lemma 3.2 all $N$ delta orientations share the single value $D_\delta$.

\section*{5. Structure of low-entropy calibrated laws}
Throughout this section and the next, $P$ is the calibrated law of an orientation $\epsilon$, $X = NP$, $D = D(P\Vert U) \le \frac{1}{60}$, and all decimal constants are safe-rounded.  Write $h_{\min} = \min_{a\ne0} h_a$.

\begin{quote}
\textbf{Lemma 5.1 (parameter floor).}  Every $a \ne 0$ satisfies
\[
|x_a| \le \sqrt{2D}
\qquad\text{and}\qquad
h_a \;\ge\; \tfrac12 \log\tfrac{1}{2D}
\;\ge\; \tfrac12\log 30 \;>\; 1.70059 .
\]
\end{quote}
\emph{Proof.}  $|x_a| = |\mathbb{E}_U[(X-1)\chi_a]| \le \mathbb{E}_U|X-1| \le \sqrt{2D}$ by Lemma 2.1.  The calibration gives $|x_a| = e^{-h_a}$, so $h_a \ge -\log\sqrt{2D} = \frac12\log\frac1{2D}$, and $D \le \frac1{60}$ gives the numerical bound ($\frac12\log 30 = 1.700598\ldots$).  $\square$

\textbf{The dip/bulk/spike decomposition.}  Partition the cube by the value of the density:

\[
\mathcal S = \{s : X(s) \le \tfrac12\} \;\text{(dips)},\quad
\mathcal T = \{s : \tfrac12 < X(s) \le 2\} \;\text{(bulk)},\quad
\mathcal S' = \{s : X(s) > 2\} \;\text{(spikes)}.
\]
For dips write $\beta_s = -\log X(s) \ge \log 2$ (the \emph{depth}; finite, since calibrated laws have full support), and define the \textbf{dip transform}

\[
\Phi(a) \;=\; \sum_{s \in \mathcal S} \beta_s\, \chi_a(s)
\qquad (a \in \mathbb{F}_2^n),
\]
with $\Phi \equiv 0$ if $\mathcal S = \varnothing$.  For spikes write $\nu_s = X(s) - 1 > 1$ and $\gamma_s = \log X(s) \in (\log 2,\, \log N]$.  Set $k = |\mathcal S|$.

\begin{quote}
\textbf{Lemma 5.2 (bookkeeping).}  With $D \le \frac1{60}$ and $\varepsilon := \sqrt{2D}$ (a scalar, not to be confused with the orientation $\epsilon$):  (i) $\;k \le \dfrac{DN}{\psi(1/2)} < 0.108630\,N$, where $\psi(\tfrac12) = \tfrac{1-\log 2}{2} = 0.153426\ldots$;  (ii) $\;\sum_{s \in \mathcal S'} \gamma_s \le \sum_{s\in\mathcal S'} \nu_s \le \tfrac{N\varepsilon}{2}$;  (iii) $\;\dfrac1N\sum_{s \in \mathcal T} |\log X(s)| \le 2\varepsilon$.
\end{quote}
\emph{Proof.}  (i) Dips have $X \le \frac12 < 1$, and $\psi$ is decreasing on $[0,1]$ (Lemma 2.2), so each dip contributes $\psi(X(s)) \ge \psi(\frac12)$ to $ND = \sum_s \psi(X(s))$, whence $k\,\psi(\tfrac12) \le ND$.  Numerically $\psi(\tfrac12) = \tfrac12\log\tfrac12 - \tfrac12 + 1 = \tfrac{1-\log2}{2}$ and $\frac{1}{60\,\psi(1/2)} = 0.1086297\ldots < 0.108630$.

(ii) $\log(1+x) \le x$ gives $\gamma_s = \log(1+\nu_s) \le \nu_s$; and $\sum_{\mathcal S'} \nu_s \le \sum_s (X(s)-1)_+ = N\,\mathbb{E}_U(X-1)_+ = \tfrac N2 \mathbb{E}_U|X-1| \le \tfrac{N\varepsilon}{2}$ by Lemma 2.1 and mean matching.

(iii) On $[\tfrac12, 2]$ we claim $|\log x| \le 2|x-1|$: for $x \in [1,2]$, $\log x \le x - 1$; for $x \in [\tfrac12, 1]$, $\log\tfrac1x \le \tfrac{1-x}{x} \le 2(1-x)$ since $x \ge \tfrac12$.  Hence $\sum_{\mathcal T}|\log X| \le 2\sum_{\mathcal T}|X-1| \le 2N\,\mathbb{E}_U|X-1| \le 2N\varepsilon$.  $\square$

\begin{quote}
\textbf{Proposition 5.3 (spectral floor and sign readout).}  Let $\varepsilon_G := \tfrac52\sqrt{2D} \le \tfrac52\sqrt{1/30} < 0.45644$, and note $h_{\min} > 2\varepsilon_G$ (indeed $1.70059 > 0.91288$).  Then for every $a \ne 0$:  (i) $\;\Big|\widehat{\log X}(a) + \dfrac{\Phi(a)}{N}\Big| \le \varepsilon_G$;  (ii) $\;|\Phi(a)| \;\ge\; N\,(h_{\min} - \varepsilon_G) \;>\; 1.24416\,N$;  (iii) $\;\epsilon_a = -\operatorname{sgn} \Phi(a)$.
\end{quote}
\emph{Proof.}  (i) Split $\widehat{\log X}(a) = \frac1N\sum_s \log X(s)\, \chi_a(s)$ over the three classes.  The dip part is exactly $-\Phi(a)/N$.  The spike part is at most $\frac1N\sum_{\mathcal S'} \gamma_s \le \frac\varepsilon2$ in absolute value by Lemma 5.2(ii), and the bulk part is at most $\frac1N\sum_{\mathcal T}|\log X| \le 2\varepsilon$ by Lemma 5.2(iii).  Total error: $\frac52\varepsilon = \varepsilon_G$.

(ii) By calibration $|\widehat{\log X}(a)| = h_a \ge h_{\min}$, so by (i), $|\Phi(a)|/N \ge h_a - \varepsilon_G \ge h_{\min} - \varepsilon_G$. Numerically $h_{\min} - \varepsilon_G > 1.70059 - 0.45644 = 1.24415$; in fact evaluating both bounds at their common worst case $D = \tfrac1{60}$ gives $h' := \tfrac12\log30 - \tfrac52\sqrt{1/30} = 1.244163\ldots > 1.24416$.  (Both $\tfrac12\log\tfrac1{2D}$ decreasing and $\tfrac52\sqrt{2D}$ increasing in $D$ make $D = \tfrac1{60}$ the simultaneous worst case.)

(iii) By (ii), $|\Phi(a)|/N \ge h_{\min} - \varepsilon_G > \varepsilon_G$, so in (i) the term $-\Phi(a)/N$ strictly dominates the error: $\widehat{\log X}(a)$ and $-\Phi(a)$ have the same (nonzero) sign.  Since $\widehat{\log X}(a) = \epsilon_a h_a$ with $h_a > 0$, $\epsilon_a = \operatorname{sgn}\widehat{\log X}(a) = -\operatorname{sgn} \Phi(a)$.  $\square$

Proposition 5.3(iii) is the pivot of the whole proof: \emph{at low entropy, the orientation is read off the deeply suppressed points.}  Note also that (ii) with $\Phi \equiv 0$ is impossible, so $\mathcal S \ne \varnothing$ automatically.

\begin{quote}
\textbf{Proposition 5.4 (dominance criterion).}  Suppose some dip $s_1 \in \mathcal S$ dominates the total depth: $\beta_{s_1} \ge \sum_{s \in \mathcal S \setminus \{s_1\}} \beta_s$. Then $\epsilon$ is the delta orientation at $s_1$.
\end{quote}
\emph{Proof.}  Write $\Phi(a) = \beta_{s_1}\chi_a(s_1) + R(a)$ with $|R(a)| \le \sum_{s \ne s_1}\beta_s \le \beta_{s_1}$.  Then $\chi_a(s_1) \Phi(a) = \beta_{s_1} + \chi_a(s_1)R(a) \ge 0$ for every $a$, and equality would force $\Phi(a) = 0$, which Proposition 5.3(ii) excludes.  Hence $\operatorname{sgn}\Phi(a) = \chi_a(s_1)$ for every $a \ne 0$, and Proposition 5.3(iii) gives $\epsilon_a = -\chi_a(s_1)$: the delta orientation at $s_1$.  $\square$

\section*{6. The deep-dip trichotomy}
\begin{quote}
\textbf{Theorem 6.1 (deep-dip trichotomy).}  Let $\epsilon$ be an orientation that is \textbf{not} a delta orientation, and suppose its calibrated law has $D \le \frac1{60}$.  Then at least three points of the cube satisfy $X(s) \le e^{-5}$, and consequently
\[
N \cdot D \;\ge\; 3\,\psi(e^{-5}) \;=\; 3\,(1 - 6e^{-5})
\;>\; 2.878716 .
\]
\end{quote}
\emph{Proof.}  Call a dip \textbf{deep} if $\beta_s \ge 5$ (i.e. $X(s) \le e^{-5}$) and \textbf{shallow} otherwise, and let $m$ be the number of deep dips.  By Lemma 5.2(i) the shallow total depth obeys

\[
\Sigma_{\mathrm{sh}} \;:=\; \sum_{\text{shallow } s} \beta_s
\;<\; 5k \;\le\; \frac{5N}{60\,\psi(1/2)} \;<\; 0.543149\,N .
\]
Recall $h' > 1.244163$ from Proposition 5.3(ii).  We rule out $m \le 2$.

\textbf{Case $m = 0$.}  Every dip is shallow, so for any $a \ne 0$ (such masks exist as $N \ge 4$),

\[
|\Phi(a)| \;\le\; \sum_{s \in \mathcal S}\beta_s \;=\;
\Sigma_{\mathrm{sh}} \;<\; 0.543149\,N \;<\; 1.24416\,N \;\le\;
|\Phi(a)|,
\]
a contradiction.  (This subsumes $k = 0$, where $\Phi \equiv 0$.)

\textbf{Case $m = 1$}, deep dip $s_1$.  For any $a \ne 0$, $|\Phi(a)| \le \beta_{s_1} + \Sigma_{\mathrm{sh}}$, so Proposition 5.3(ii) forces

\[
\beta_{s_1} \;\ge\; N h' - \Sigma_{\mathrm{sh}}
\;>\; (1.244163 - 0.543149)\,N \;=\; 0.701014\,N
\;>\; 0.543149\,N \;>\; \Sigma_{\mathrm{sh}}
\;=\; \sum_{s \ne s_1}\beta_s ,
\]
so $s_1$ dominates and Proposition 5.4 says $\epsilon$ is a delta orientation — contradiction.

\textbf{Case $m = 2$}, deep dips $s_1, s_2$ with (without loss of generality) $\beta_{s_1} \ge \beta_{s_2}$.  Since $t := s_1 + s_2 \ne 0$, the set $\{a : \chi_a(t) = -1\}$ has exactly $N/2$ elements and does not contain $a = 0$; pick such an $a$.  For it, $\chi_a(s_2) = -\chi_a(s_1)$, so

\[
|\Phi(a)| \;\le\; |\beta_{s_1} - \beta_{s_2}| + \Sigma_{\mathrm{sh}},
\qquad\text{whence}\qquad
\beta_{s_1} - \beta_{s_2} \;\ge\; Nh' - \Sigma_{\mathrm{sh}} \;>\;
0.701014\,N .
\]
Then

\[
\beta_{s_1} - \Big(\beta_{s_2} + \Sigma_{\mathrm{sh}}\Big)
\;>\; 0.701014\,N - 0.543149\,N \;=\; 0.157865\,N \;>\; 0,
\]
so again $s_1$ dominates the total depth and Proposition 5.4 gives a contradiction.

\textbf{Conclusion.}  $m \ge 3$.  Deep dips have $X(s) \le e^{-5} < 1$ and $\psi$ is decreasing on $[0,1]$, so each contributes $\psi(X(s)) \ge \psi(e^{-5}) = e^{-5}(-5) - e^{-5} + 1 = 1 - 6e^{-5}$ to $ND = \sum_s\psi(X(s))$ (Lemma 2.2), and all terms are nonnegative.  Hence $ND \ge 3(1 - 6e^{-5}) = 2.8787169\ldots > 2.878716$.  $\square$

\section*{7. Proof of the main theorem for $n \ge 6$}
\begin{quote}
\textbf{Theorem 7.1.}  Let $n \ge 6$ and let $\epsilon$ be any orientation that is not a delta orientation.  Then $\widehat m(\epsilon) > D_\delta$.
\end{quote}
\emph{Proof.}  Write $D = \widehat m(\epsilon)$ and split into two cases.

\textbf{Case $D > \frac1{60}$.}  Since $n \ge 6$, $N - 1 \ge 63$, so Lemma 4.2 gives

\[
D_\delta \;<\; \frac1{N-1} \;\le\; \frac1{63} \;<\; \frac1{60}
\;<\; D .
\]
\textbf{Case $D \le \frac1{60}$.}  Theorem 6.1 applies and gives $D \ge 2.878716/N$.  Since $2.878716\,(N-1) > N$ for every $N \ge 2$ (equivalently $N > \tfrac{2.878716}{1.878716} = 1.532\ldots$),

\[
D \;\ge\; \frac{2.878716}{N} \;>\; \frac{1}{N-1} \;>\; D_\delta .
\]
(For $n = 6, 7$ this case is in fact vacuous — the conclusion of Theorem 6.1 contradicts $D \le \frac1{60}$ when $N < 60 \cdot 2.878716$ — which is a legitimate outcome of the case split: those $n$ are carried entirely by the first case.)

Both cases give strict inequality.  $\square$

\emph{Proof of Theorem 1.2 and Corollary 1.3.}  For $n \ge 6$: Theorem 7.1 gives strict inequality for every non-delta orientation, and all $N$ delta orientations share the value $D_\delta$ by Lemma 3.2; the equality set is therefore exactly the delta family, and the quantitative floor of Corollary 1.3 is the case split just proved ($D > \frac1{60}$ gives $ND > N/60$; otherwise $ND \ge 2.878716$), together with $N D_\delta < N/(N-1) \le 64/63$ from Lemma 4.2.  For $2 \le n \le 5$ the statement is Theorem 8.3 below (certified computation).  $\square$

\section*{8. The certified computation for $2 \le n \le 5$}
This section proves the main theorem in the range the analytic argument does not reach.  The computation is organized so that every reported number carries an a posteriori error certificate resting on the two lemmas below; the code and canonical outputs ship with this repository (Section 10), and the trust base is stated in full in Section 8.5.

\subsection*{8.1 Certification lemmas}
The solver minimizes the strictly convex $G_\epsilon$ of Theorem 3.1 by damped Newton iteration.  Certification does not trust the iteration: it evaluates, at the numeric output $\tilde\ell$, the gradient residual and a curvature floor, and converts them into a rigorous error radius.

\begin{quote}
\textbf{Lemma 8.1 (a posteriori radius).}  Fix an orientation $\epsilon$ and any $\tilde\ell \in \mathbb{R}^{N-1}$.  Let $\rho = \Vert\nabla G_\epsilon(\tilde\ell)\Vert_2$ and $\lambda = \min_a e^{-\epsilon_a\tilde\ell_a}$.  If $r_0 := 2e\rho/\lambda \le 1$, then the true minimizer satisfies $\Vert\ell^\ast - \tilde\ell\Vert_2 \le r_0$.
\end{quote}
\emph{Proof.}  If $\rho = 0$ then $\tilde\ell = \ell^\ast$ by strict convexity.  Otherwise, suppose $T := \Vert\ell^\ast - \tilde\ell\Vert_2 > r_0$ and let $v = (\ell^\ast - \tilde\ell)/T$.  The function $g(t) = \langle \nabla G_\epsilon(\tilde\ell + tv),\, v\rangle$ is nondecreasing (convexity) with $g(T) = 0$ and $g(0) \ge -\rho$ (Cauchy–Schwarz).  On the segment $t \in [0, r_0]$, each coordinate moves by at most $|t v_a| \le t \le 1$, so the barrier part of the Hessian obeys $e^{-\epsilon_a(\tilde\ell_a + t v_a)} \ge \lambda e^{-1}$, and $\nabla^2 G_\epsilon \succeq \operatorname{diag}(e^{-\epsilon_a \ell_a}) \succeq \lambda e^{-1} I$ there (the covariance part is positive semidefinite).  Hence

\[
g(r_0) \;\ge\; g(0) + r_0\,\lambda e^{-1}
\;\ge\; -\rho + \frac{2e\rho}{\lambda}\cdot\frac{\lambda}{e}
\;=\; \rho \;>\; 0 ,
\]
contradicting $g$ nondecreasing with $g(T) = 0$ at $T > r_0$. $\square$

\begin{quote}
\textbf{Lemma 8.2 (entropy transfer).}  Let $\mathcal D(\ell) := D(P_\ell \Vert U) = \langle \ell, x(\ell)\rangle - F(\ell)$.  Then $\nabla\mathcal D(\ell) = \nabla^2 F(\ell)\,\ell$ and $\Vert\nabla^2 F(\ell)\Vert_{\mathrm{op}} \le N - 1$; consequently, in the setting of Lemma 8.1,
\[
\big|\mathcal D(\ell^\ast) - \mathcal D(\tilde\ell)\big|
\;\le\; (N-1)\,\big(\Vert\tilde\ell\Vert_2 + r_0\big)\, r_0 .
\]
\end{quote}
\emph{Proof.}  $\log X_\ell = \sum_a \ell_a\chi_a - F(\ell)$, so $\mathcal D(\ell) = \mathbb{E}_{P_\ell}[\log X_\ell] = \langle\ell, x(\ell)\rangle - F(\ell)$ and $\nabla \mathcal D = x + \nabla^2 F\,\ell - x = \nabla^2 F\,\ell$. The Hessian $\nabla^2 F = \operatorname{Cov}_{P_\ell}(\chi) \preceq \mathbb{E}_{P_\ell}[\chi\chi^{\mathsf T}]$ in the positive semidefinite order (subtracting the rank-one $xx^{\mathsf T}$), and a positive semidefinite matrix has operator norm at most its trace: $\operatorname{tr}\,\mathbb{E}[\chi\chi^{\mathsf T}] = \sum_{a\ne0}\mathbb{E}[\chi_a^2] = N-1$.  The transfer bound is the mean value inequality along the segment, whose points have norm at most $\Vert\tilde\ell\Vert_2 + r_0$.  $\square$

\textbf{Delta reference values.}  $u_\star$ is bracketed by sign-of-$p$ bisection (Proposition 4.1: $p$ strictly increasing, $p(0) < 0 < p(\tfrac1{N-1})$) to width $2^{-300}$ at 60 significant digits, and $D_\delta = \frac1N[A\log A + (N-1)B\log B]$ is evaluated at both bracket ends.  On the bracket, $D_\delta(u)$ is strictly increasing, $\frac{d}{du}D_\delta(u) = \frac{N-1}{N}\log\frac{B}{A} > 0$ — so the two endpoint evaluations enclose the true value up to the arithmetic error of the evaluations themselves.

\subsection*{8.2 Exhaustive certification for $n = 2, 3, 4$}
For each of the $2^{N-1}$ orientations ($8$, $128$, $32768$): solve by damped Newton in double precision; re-evaluate $\rho$, $\lambda$, and $\mathcal D(\tilde\ell)$ at the exact binary-rational point $\tilde\ell$ in 30-significant-digit arithmetic; apply Lemmas 8.1 and 8.2; and compare the certified enclosure of $\mathcal D(\ell^\ast)$ against the certified delta enclosure.  Outcomes (canonical outputs in \texttt{receipts/}):

\begin{center}
\begin{adjustbox}{max width=\textwidth}
\begin{tabular}{llllll}
$n$ & orientations & failures & worst certified non-delta margin & max $r_0$ & max entropy error \\
\hline
2 & 8 & 0 & 0.193236 & $5.8\times10^{-14}$ & $3.4\times10^{-13}$ \\
3 & 128 & 0 & 0.462239 & $2.8\times10^{-14}$ & $5.5\times10^{-13}$ \\
4 & 32768 & 0 & 0.263509 & $2.3\times10^{-11}$ & $1.7\times10^{-9}$ \\
\end{tabular}
\end{adjustbox}
\end{center}
Each margin is (certified lower bound on the non-delta value) minus (certified upper bound on $D_\delta$); every delta orientation's enclosure matched the closed-form reference.  The certified errors are 8–12 orders of magnitude below the margins.

\subsection*{8.3 Orbit reduction and certification for $n = 5$}
At $n = 5$ there are $2^{31}$ orientations; exhaustive certification is avoided using the symmetry group $\Gamma_5 = \langle\text{ translations}, GL(5,2)\rangle$, under which $\widehat m$ is invariant (Lemmas 3.2 and 3.3) and the delta family forms one orbit.

\emph{Generators.}  Translations by the $n$ basis vectors realize all translations.  For $GL(5,2)$ we use the coordinate cycle, the transposition of the first two coordinates, and the transvection $e_2 \mapsto e_1 + e_2$: transpositions plus the cycle generate all coordinate permutations, and conjugating the transvection by permutations yields every elementary transvection $e_j \mapsto e_i + e_j$; elementary transvections generate $GL(n,2)$ (row reduction over $\mathbb{F}_2$).

\emph{Enumeration.}  A breadth-first search over the full state space $\{\pm1\}^{31}$ (bitmask representation, $2^{31}$ states, a visited bitmap) partitions it into orbits of the generated group.  The partition is complete \emph{by construction}; as a certificate, the orbit sizes are summed and checked to equal $2^{31}$ exactly.  Result: \textbf{176 orbits}.  Two independent cross-checks: (i) at $n = 2, 3, 4$ the same BFS code yields $2$, $4$, $14$ orbits, matching the exhaustive enumerations of Section 8.2; (ii) an independent Burnside-lemma computation (receipt \texttt{w5}): for each $M \in GL(n,2)$ the elements $(t, M)$ of $\Gamma_n$ jointly fix $2^{c(M) + n - r(M)}$ orientations, where $c(M)$ is the number of cycles of $a \mapsto M^{\mathsf T}a$ on nonzero masks and $r(M)$ the $\mathbb{F}_2$-rank of the cycle-XOR values (a $\sigma$ fixed by $(t,M)$ needs $\chi_{X_j}(t) = 1$ for each cycle-XOR $X_j$, leaving $2^{c(M)}$ free signs); averaging over all $9{,}999{,}360$ elements of $GL(5,2)$ gives exactly $176$, with the divisibility check exact, and the same program reproduces $2, 4, 14$ at $n = 2, 3, 4$.  (Even if the generators had generated a proper subgroup, the enumeration would only \emph{refine} the true orbit partition and the certification below would remain valid — the Burnside match shows it is in fact the full partition.)

\emph{Certification.}  Each of the 176 representatives is solved with a double-precision warm start plus Newton refinement at 40 significant digits, then certified via Lemmas 8.1–8.2.  Outcomes: zero failures; the delta orbit (size 32) is uniquely minimal with

\[
D_\delta(n{=}5) = 0.0317486983145803\ldots,
\]
matching the closed-form reference; the worst certified non-delta margin is $0.138872$ (attained by the orbit of size 4960 described in Section 9), with certified entropy errors below $10^{-32}$ on every orbit (the receipt prints the global maximum; the five tabulated orbits sit below $10^{-33}$).  The five lowest orbit values:

\begin{center}
\begin{adjustbox}{max width=\textwidth}
\begin{tabular}{lll}
orbit representative & orbit size & certified $\widehat m$ \\
\hline
delta orbit & 32 & 0.031748698 \\
three-dip family (Section 9) & 4960 & 0.170621384 \\
— & 27776 & 0.234706844 \\
$j{=}3$ family (Section 9) & 34720 & 0.243617531 \\
— & 317440 & 0.405832442 \\
\end{tabular}
\end{adjustbox}
\end{center}
\begin{quote}
\textbf{Theorem 8.3.}  For $2 \le n \le 5$, every non-delta orientation satisfies $\widehat m(\epsilon) > D_\delta$, with certified margin at least $0.138$; the equality set of Theorem 1.2 is exactly the delta family.  \emph{(Computer-assisted; trust base in Section 8.5.)}
\end{quote}
\subsection*{8.4 Independent replication}
During development, all computations were additionally replicated on independently written code paths (separate solver implementations, including a log-domain variant): the $n = 4$ exhaustive scan reproduced every value; the 176 orbit representatives reproduced to $\le 4\times10^{-10}$; and an adversarial search campaign (about $5.5\times10^4$ solves at $4 \le n \le 13$, including designed attacks on the trichotomy of Theorem 6.1) found no violation of any inequality asserted in this paper — in particular, no calibrated non-delta law with fewer than three deep dips was found at any entropy level.  These replications are \emph{not} shipped with this repository and none of the certificates above rely on them; the shipped, rerunnable evidence is receipts \texttt{w1}–\texttt{w5}.

\subsection*{8.5 Trust base}
The $n \le 5$ leg of Theorem 1.2 is computer-assisted in the following precise sense.  Arithmetic is faithful round-to-nearest multiple-precision floating point (mpmath [5], 30–60 significant digits), not directed-rounding interval arithmetic; the certified margins ($\ge 0.138$) exceed the certified error bounds ($\le 1.7\times10^{-9}$) by 8 orders of magnitude and the residual arithmetic uncertainty of the evaluations themselves by far more, and all values were independently replicated (Section 8.4), but a fully formal verification would substitute outward-rounded interval arithmetic throughout — a mechanical upgrade that we have not performed.  The two C programs (\texttt{w2}, \texttt{w5}) use exact 64-bit integer arithmetic only — no floating point — so they require no rounding analysis; they compile with any C99 compiler (\texttt{cc -O2}; the canonical outputs were produced with Apple clang).  The remaining trust items are: correctness of Lemmas 8.1–8.2 (proved above), completeness of the orbit partition (checksum plus the independent Burnside count of receipt \texttt{w5}), and the invariance lemmas 3.2–3.3 (proved above).  The analytic part of the paper ($n \ge 6$) uses no computation.

\section*{9. Sharpness: the three-dip family and the runner-up constant}
The proof of Theorem 6.1 shows any cheap non-delta orientation must extinguish at least three points.  There is a canonical family that extinguishes \emph{exactly} three, and it appears to be the true runner-up. Fix two distinct nonzero masks $t_1, t_2$ (say the first two coordinate masks) and define

\[
\epsilon_a = \begin{cases}
+1, & \chi_a(t_1) = \chi_a(t_2) = -1
\text{ (both defining bits of } a \text{ set)},\\
-1, & \text{otherwise.}
\end{cases}
\]
Equivalently: start from the delta at $0$ and flip the signs on the quarter of the masks containing both distinguished bits.  Its calibrated law (computed by the certified solver; this is the size-4960 orbit of Section 8.3 at $n=5$, and the runner-up in the exhaustive scans at $n = 3, 4$) extinguishes the three points $\{0, e_1, e_2\}$ of the affine plane $\{0, e_1, e_2, e_1{+}e_2\}$, boosts the fourth plane point to $X \to 4$, and is asymptotically uniform elsewhere.  Its entropy gap satisfies, numerically,

\begin{center}
\begin{adjustbox}{max width=\textwidth}
\begin{tabular}{lllllll}
$n$ & 3 & 4 & 5 & 6 & 7 & 8 \\
\hline
$N \cdot \widehat m$ & 4.7662 & 5.2488 & 5.4599 & 5.5233 & 5.5397 & 5.5438 \\
\end{tabular}
\end{adjustbox}
\end{center}
converging to $4\log4 = 5.545177\ldots$, which is exactly $3\,\psi(0) + \psi(4)$ — three extinguished points at limit cost $1$ each, plus one point at density $4$ (the value forced by conservation of mass within the plane).  The $n \le 5$ entries are certified (Sections 8.2–8.3); the trend in $n$ is monotone increasing.  The generalization that flips signs on the masks containing $j$ fixed bits extinguishes $2^{j-1}+1$ points and satisfies $N\widehat m \to 2^j\log\frac{2^j}{2^{j-1}-1}$ numerically ($7.84663\ldots$ at $j=3$), increasing in $j$.  (The case $j = 1$ is not degenerate but exact: flipping the masks that contain one fixed bit turns the delta at $0$ into the delta at $e_1$, whose constant is $N D_\delta \to 1$; the displayed formulas apply for $j \ge 2$.)

\begin{quote}
\textbf{Conjecture 9.1 (sharp runner-up).}  $\displaystyle \liminf_{n\to\infty}\ \min_{\epsilon \text{ non-delta}}\ N\cdot\widehat m(\epsilon) \;=\; 4\log 4 ,$ attained along the three-dip family.
\end{quote}
The gap between the proved floor $2.878716$ (Theorem 6.1) and the conjectured $5.545$ is the price of the trichotomy's generosity to shallow dips; closing it would require classifying the limit configurations of calibrated laws at entropy scale $1/N$, which we leave open.  We note that Theorem 1.2 itself needed only the floor to beat $N D_\delta < 64/63$ — the margin is $2.8\times$, which is why the theorem closes while the sharp constant remains open.

\section*{10. Code and reproducibility}
All code is in \texttt{code/} and all canonical outputs in \texttt{receipts/} of this repository.

\begin{itemize}
\item \texttt{w1\_cert\_small\_n.py} — exhaustive certified verification for $n = 2, 3, 4$ (Section 8.2): double-precision damped Newton on $G_\epsilon$, 30-digit re-evaluation at the exact output point, Lemmas 8.1–8.2, delta bracketing at 60 digits.  Receipt: \texttt{w1\_cert\_small\_n.out} ($n=4$ takes roughly an hour on a laptop).
\item \texttt{w2\_orbit\_bfs.c} — the $\Gamma_n$-orbit enumeration by BFS with completeness checksum (Section 8.3), for $n \le 5$; prints the summary line (orbit count and checksum) to stdout and writes one line per orbit (minimal representative, size) to \texttt{orbits\_n<n>.txt}. Receipts: \texttt{w2\_orbit\_bfs.out} (the four summary lines, counts $2/4/14/176$) and \texttt{w2\_orbits\_n5.txt} (the 176-line orbit list).
\item \texttt{w3\_cert\_n5.py} — certified solves of the 176 representatives (Section 8.3), 40-digit refinement plus certificates; reads the orbit list from \texttt{receipts/} by a script-relative path (runs from any working directory) and writes a per-orbit array to \texttt{/tmp}.  Receipt: \texttt{w3\_cert\_n5.out}.
\item \texttt{w5\_burnside.c} — the independent Burnside-lemma orbit count of Section 8.3 (all invertible matrices enumerated; exact divisibility checked).  Receipt: \texttt{w5\_burnside.out}.
\item \texttt{w4\_family\_receipt.py} — the three-dip and $j{=}3$ families at $n = 9, 10, 11$: verifies in log-domain every inequality of Sections 5–6 on genuinely low-entropy laws (parameter floor, spectral floor, sign readout, deep-dip count, the trichotomy conclusion), and prints the Section 9 table.  Receipt: \texttt{w4\_family\_receipt.out}.  \emph{Numerical note:} depths reach $\beta \sim N\log(N/4)$, far below double precision underflow ($X \sim e^{-5700}$ at $n=10$); all depth-dependent quantities must be computed in log-domain, i.e. $\log X = \log N + \text{field} - \operatorname{logsumexp}$, never by exponentiating.
\end{itemize}
\section*{Declarations}
No funding was received for this work.  The author has no competing interests.  All code and canonical outputs needed to verify the computer-assisted claims ship with this paper (Section 10).

\section*{References}
[1] T. M. Cover and J. A. Thomas, \emph{Elements of Information Theory}, 2nd ed., Wiley, 2006.  (Pinsker's inequality: Lemma 11.6.1.)

[2] R. O'Donnell, \emph{Analysis of Boolean Functions}, Cambridge University Press, 2014.

[3] T. Erdélyi, Recent progress in the study of polynomials with constrained coefficients, in \emph{Trigonometric Sums and Their Applications}, Springer, 2020, pp. 29–69.

[4] M. J. Wainwright and M. I. Jordan, Graphical models, exponential families, and variational inference, \emph{Foundations and Trends in Machine Learning} 1 (2008), 1–305.

[5] F. Johansson et al., \emph{mpmath: a Python library for arbitrary-precision floating-point arithmetic}, version 1.3.0 (2023), https://mpmath.org/.

[6] M. Mézard and A. Montanari, \emph{Information, Physics, and Computation}, Oxford University Press, 2009.

\end{document}
